\begin{corollary}[7.5.2]
\label{IV.2.s7.5.2-fr}
Sean $A$ un anillo local noetheriano, $\mathfrak m$ el ideal maximal de $A$,
$k=A/\mathfrak m$ su cuerpo residual, $B$ un anillo local noetheriano
y $\varphi:A\to B$ un homomorfismo local tal que:
\begin{enumerate}
\item[\emph{(i)}] El cuerpo residual de $B$ es una extensión finita de $k$.
\item[\emph{(ii)}] $B$ es un $A$-módulo plano.
\end{enumerate}
Sea, por otra parte, $\boldsymbol R(C)$ una propiedad que satisface las
condiciones $(\mathrm R_{\mathrm{II}})$ \textup{(7.3.10)},
$(\mathrm R_{\mathrm{III}})$ \textup{(7.3.18)} y $(\mathrm R_{\mathrm{IV}})$
\textup{(7.5.1)}. Por último, sea $\boldsymbol R'(C)$ una propiedad que satisface
la condición siguiente:
\begin{quote}
$(\mathrm R_{\mathrm V})$ Si $C$ y $D$ son dos anillos locales
noetherianos y $\psi:C\to D$ es un homomorfismo local tal que
\[
g={}^a\psi:\operatorname{Spec}(D)\to\operatorname{Spec}(C)
\]
es un $\boldsymbol P$-morfismo, entonces la propiedad
$\boldsymbol R'(C)$ implica $\boldsymbol R(D)$.
\end{quote}
Supongamos que el morfismo canónico
$\operatorname{Spec}(\widehat A)\to\operatorname{Spec}(A)$ es un
$\boldsymbol P'$-morfismo, donde $\boldsymbol P'$ se define a partir
de $\boldsymbol R'$ como $\boldsymbol P$ a partir de $\boldsymbol R$
en \textup{(7.5.0)}.

En estas condiciones, sea
$f={}^a\varphi:\operatorname{Spec}(B)\to\operatorname{Spec}(A)$ y
supongamos que la propiedad
$\boldsymbol P(\operatorname{Spec}(\widehat B
\otimes_{\widehat A}k),k)$
es verdadera. Entonces, para todo $x\in\operatorname{Spec}(A)$, la propiedad
$\boldsymbol P(f^{-1}(x),\boldsymbol{k}(x))$ es verdadera \textup{(es decir,} $f$ \textup{es un}
$\boldsymbol P$-morfismo).
\end{corollary}

\begin{proof}
Procedamos de nuevo en dos etapas:

I) \emph{Reducción al estudio de los anillos locales de la fibra genérica.}
--- Apliquemos de nuevo el lema (7.3.16.2). La única diferencia con el
razonamiento de I) en (7.5.1) es que aquí $A'$ es un anillo \emph{semilocal}
íntegro, pero no necesariamente local; todo ideal maximal
$\mathfrak n'$ de $B'$ está entonces sobre un ideal maximal
$\mathfrak m'$ de $A'$. Como todo anillo local de
$\operatorname{Spec}(B'\otimes_{A'}K')$ es un anillo local de
$\operatorname{Spec}(B')$ y, por tanto, de uno de los
$\operatorname{Spec}(B'_{\mathfrak n'})$ (sobre el punto genérico de
$\operatorname{Spec}(A'_{\mathfrak m'}))$, la cuestión se reduce a
demostrar que los anillos locales de
$\operatorname{Spec}(B'_{\mathfrak n'}
\otimes_{A'_{\mathfrak m'}}K')$
poseen la propiedad $\boldsymbol R$. Ahora bien, la definición de
$\boldsymbol P'$ implica que, si se cumple $\boldsymbol P'(Z,k)$,
también se cumple $\boldsymbol P'(Z\otimes_k k',k')$ para toda extensión
finita $k'$ de $k$. El mismo razonamiento que en (7.3.7) demuestra que
las fibras formales de $A'_{\mathfrak m'}$ poseen la propiedad
$\boldsymbol P'$. Como en la parte I) de la demostración de (7.5.1),
se ve que $A'_{\mathfrak m'}$ y $B'_{\mathfrak n'}$ satisfacen las
condiciones (i) y (ii) de (7.5.2). Por otra parte, si $k'$ es el cuerpo
residual de $A'_{\mathfrak m'}$, la propiedad
$\boldsymbol P(\operatorname{Spec}(\widehat{B}'_{\mathfrak n'}
\otimes_{\widehat{A}'_{\mathfrak m'}}k'),k')$
es verdadera: en efecto, el completado de $A'_{\mathfrak m'}$ es uno de los anillos
locales componentes de $\widehat{A}'=\widehat A\otimes_A A'$, y el
completado de $B'_{\mathfrak n'}$ es uno de los anillos locales
componentes de
$\widehat{B}'=\widehat B\otimes_B B'
=\widehat B\otimes_A A'
=\widehat B\otimes_{\widehat A}\widehat{A}'$;
el razonamiento de I) en (7.5.1) demuestra, por tanto, nuestra afirmación.

Es posible, pues, limitarse al caso en que $A$ es además \emph{íntegro}, y a
demostrar que, si $K$ es el cuerpo de fracciones de $A$, \emph{los anillos
locales de} $\operatorname{Spec}(B\otimes_A K)$ \emph{poseen la propiedad}
$\boldsymbol R$.

II) \emph{Reducción al caso en que $A$ y $B$ son completos.} --- Sean
$\xi$ el punto genérico de $\operatorname{Spec}(A)$ e $y$ un punto
cualquiera de $f^{-1}(\xi)$. Se trata de demostrar que se cumple
$\boldsymbol R(\mathscr O_y)$. Teniendo en cuenta la condición
$(\mathrm R_{\mathrm{II}})$, bastará ver que existe en
$\operatorname{Spec}(\widehat B)$ un punto $z$ sobre $y$ para el que
se cumple $\boldsymbol R(\mathscr O_z)$. Ahora bien, de la hipótesis
sobre $\operatorname{Spec}(\widehat B\otimes_{\widehat A}k)$ y de
(7.5.1) se deduce que el morfismo
${}^a\widehat\varphi:\operatorname{Spec}(\widehat B)
\to\operatorname{Spec}(\widehat A)$
es un $\boldsymbol P$-morfismo. Por otra parte, para todo punto
$z\in\operatorname{Spec}(\widehat B)$ sobre $y$ (existe tal punto,
pues $\widehat B$ es un $B$-módulo fielmente plano), la imagen $x$ de
$z$ en $\operatorname{Spec}(\widehat A)$ pertenece a la fibra
$h^{-1}(\xi)$ del morfismo canónico
$h:\operatorname{Spec}(\widehat A)\to\operatorname{Spec}(A)$. En
virtud de la hipótesis sobre $A$, se cumple la propiedad
$\boldsymbol R'(\mathscr O_x)$; por tanto, en virtud de
$(\mathrm R_{\mathrm V})$, se cumple $\boldsymbol R(\mathscr O_z)$.
Queda demostrado.
\end{proof}
